\documentclass[12pt]{article}
% Préambule commun à tous les documents extraits de « La Revue CAPES »
\usepackage[T1]{fontenc}
\usepackage[utf8]{inputenc}
\usepackage{lmodern}
\usepackage{amsmath,amssymb,amsthm,amsfonts,mathrsfs}
\usepackage{setspace}
\usepackage{listings}
\usepackage{pgfplots}
\pgfplotsset{compat=1.18}
\usepackage{longtable}
\usepackage{adjustbox}
\usepackage{lettrine}
\usepackage{tkz-tab}
\usepackage{changepage}
\usepackage{pdflscape}
\usepackage{graphicx}
\usepackage{tikz}
\usepackage{xcolor}
\usepackage[a4paper,top=2cm,bottom=2.5cm,left=2.5cm,right=2cm]{geometry}
\usepackage{fancyhdr}
\usepackage{draftwatermark}
\SetWatermarkText{La RevueCapes}
\SetWatermarkLightness{0.9}
\SetWatermarkScale{2}
\usepackage{hyperref}
\hypersetup{colorlinks=true,linkcolor=blue!50!black,urlcolor=blue!50!black}

\definecolor{revue}{RGB}{20,60,120}

\pagestyle{fancy}
\fancyhf{}
\renewcommand{\headrulewidth}{0.4pt}
\setlength{\headheight}{15pt}

% \entete{Matière}{Titre}{Type de document}
\newcommand{\entete}[3]{%
  \fancyhead[L]{\small\textsc{La Revue CAPES} -- #1}%
  \fancyhead[R]{\small #2}%
  \fancyfoot[C]{\thepage}%
  \thispagestyle{plain}%
  \begin{center}
    {\color{revue}\rule{\textwidth}{1.2pt}}\\[4pt]
    {\small\textsc{Concours directs CAP-CEG / CAPES -- Burkina Faso}}\\[6pt]
    {\LARGE\bfseries\color{revue} #2}\\[6pt]
    {\large\itshape #1 -- #3}\\[2pt]
    {\color{revue}\rule{\textwidth}{1.2pt}}
  \end{center}
  \vspace{0.5cm}%
}

% Encadré pour les remarques ajoutées dans les corrigés
\newcommand{\remarque}[1]{\par\medskip\noindent\fcolorbox{revue}{revue!6}{\parbox{\dimexpr\linewidth-2\fboxsep-2\fboxrule}{\textbf{\color{revue}Remarque.} #1}}\par\medskip}


\begin{document}
\entete{Culture générale}{Corrigé du sujet type de dissertation n°6}{Correction}
\begin{spacing}{1.5}

\textbf{Introduction}

Albert Jacquard, biologiste et humaniste, met en avant une vision globale et universelle de la solidarité. Cette citation, "Désormais la solidarité la plus nécessaire est celle de l'ensemble des habitants de la terre", insiste sur l'importance de l'interdépendance entre les êtres humains dans un monde de plus en plus globalisé. Dans un contexte marqué par des crises multiples — environnementales, économiques, sanitaires, et sociales — cette affirmation nous invite à réfléchir à l'urgence d'une solidarité mondiale. Nous analyserons cette idée à travers trois axes : la nécessité d'une solidarité universelle, les défis à surmonter pour l'atteindre, et les limites éventuelles de cette perspective.

1. La nécessité d'une solidarité universelle

1.1. L'interconnexion des défis mondiaux
Les crises modernes, telles que le réchauffement climatique, les pandémies (comme le COVID-19), et les inégalités sociales, affectent tous les habitants de la planète sans distinction de frontières.

Exemple : Le changement climatique, responsable de catastrophes naturelles croissantes, exige une action collective. Les pays pollueurs doivent réduire leurs émissions pour éviter des conséquences désastreuses pour toute l'humanité.

Analyse : Ces enjeux ne peuvent être résolus qu'avec une solidarité globale, car ils dépassent les capacités individuelles des nations.

1.2. Une humanité interdépendante
Les progrès technologiques et économiques ont intensifié l'interdépendance entre les peuples.

Exemple : La chaîne d'approvisionnement mondiale illustre que le bien-être d'un pays dépend de la stabilité des autres. Une crise en Asie peut rapidement affecter l'Europe ou l'Afrique.

Analyse : La solidarité n'est plus une option morale, mais une nécessité fonctionnelle pour maintenir un équilibre global.

1.3. Un impératif éthique
La solidarité mondiale repose aussi sur une reconnaissance de l'humanité partagée.

Exemple : La Déclaration universelle des droits de l'homme affirme que tous les humains ont droit à une vie digne.

Analyse : L'idéal de solidarité repose sur la conviction que nous sommes coresponsables du bien-être collectif.

2. Les défis de la solidarité mondiale

2.1. Les inégalités entre les nations
Les écarts économiques et sociaux entre les pays riches et pauvres rendent la solidarité difficile à mettre en œuvre.

Exemple : Lors de la pandémie de COVID-19, les pays riches ont monopolisé les vaccins, privant les pays pauvres d'un accès équitable.

Analyse : Tant que les intérêts économiques domineront, la solidarité mondiale restera inégale et insuffisante.

2.2. Les conflits d'intérêts nationaux
Chaque pays tend à privilégier ses propres intérêts, même au détriment du bien commun.

Exemple : Les accords climatiques, comme celui de Paris, peinent à être respectés en raison des réticences de certains États à sacrifier leur croissance économique.

Analyse : La solidarité nécessite des compromis, mais les rivalités géopolitiques entravent souvent cette coopération.

2.3. Le manque de gouvernance mondiale efficace

L'absence d'une autorité mondiale capable d'imposer des décisions justes et équitables limite les efforts collectifs.

Exemple : Les Nations Unies, bien qu'influentes, n'ont pas les moyens d'imposer des solutions globales.

Analyse : Sans structures solides, la solidarité mondiale repose sur la bonne volonté des États, qui n'est pas toujours garantie.

3. Les limites et conditions de la solidarité universelle

3.1. Les différences culturelles et idéologiques

Les visions divergentes du monde freinent l'adoption d'un consensus global.

Exemple : Les notions de solidarité varient entre les cultures occidentales, axées sur les droits individuels, et les sociétés communautaires, privilégiant le collectif.

Analyse : Une solidarité mondiale doit respecter ces diversités tout en établissant des principes communs.

3.2. La complexité des responsabilités partagées

Tous les pays n'ont pas les mêmes responsabilités dans les crises actuelles.

Exemple : Les nations industrialisées sont les principales responsables du changement climatique, mais les pays en développement en subissent les conséquences.

Analyse : La solidarité mondiale doit inclure une justice climatique, où les plus responsables contribuent davantage.

3.3. La méfiance et les inégalités historiques

L'histoire des colonisations et des injustices passées alimente la méfiance entre nations.

Exemple : Certains pays africains voient les initiatives solidaires des nations riches comme une forme de néocolonialisme.

Analyse : La solidarité doit être dénuée de paternalisme et fondée sur un respect mutuel.

\textbf{Conclusion}

Albert Jacquard a raison de souligner que la solidarité entre tous les habitants de la Terre est aujourd'hui indispensable. Face à des défis mondiaux comme le climat, la santé et les inégalités, l'humanité doit unir ses forces pour assurer un avenir durable. Cependant, la mise en œuvre de cette solidarité se heurte à des obstacles politiques, économiques, et culturels. Construire une solidarité universelle nécessite des efforts collectifs, mais aussi une justice équitable, afin de surmonter les inégalités historiques et les divergences idéologiques. C'est un défi ambitieux, mais essentiel pour l'avenir de notre planète et de l'humanité.




















\quad
\quad

\end{spacing}
\end{document}
